\documentclass[12pt]{lettre} 
\usepackage[utf8]{inputenc} 
\usepackage[T1]{fontenc} 
\usepackage{lmodern} 
\usepackage{eurosym} 
\usepackage{numprint} 

% Lettre de demande de modification du systeme de videoprotection d'Evres.
% Champs a completer : [localisations sites D et E], [faits recents], [date].

\begin{document}

\begin{letter}{Monsieur le Préfet de la Meuse\\Préfecture de la Meuse\\Bureau de l'ordre public et de la sécurité intérieure\\40 rue du Bourg\\55012 BAR LE DUC CEDEX}


\address{Clément \textsc{Fevez}\\ Maire de Evres\\Place de la Halle\\55250 EVRES}
\lieu{Evres}
\telephone{03 29 92 05 97}
\email{evres@orange.fr}
\nofax


\def\concname{Objet :~}
\conc{Demande de modification de l'autorisation d'installation du système de vidéoprotection d'Evres (arrêté n° 2026-387 du 10 mars 2026)}

\opening{Monsieur le Préfet,}
J'ai l'honneur de solliciter la modification de l'autorisation d'installation du système de vidéoprotection de la commune d'Evres, que vous m'avez accordée par votre arrêté n° 2026-387 du 10 mars 2026 sur la demande enregistrée sous le n° 20240253, et qui demeure en cours de validité.

Cette modification consiste :

\begin{itemize}
\item en l'installation de deux caméras nouvelles (caméras 7 et 8) sur un nouveau site D, couvrant [zone / voie non couverte à ce jour, à préciser] ;
\item en le déplacement de la caméra 1, aujourd'hui installée sur le site A où son champ de vision est redondant avec celui de la caméra 2, vers un nouveau site E couvrant [zone / voie non couverte à ce jour, à préciser].
\end{itemize}

Le dispositif porterait ainsi à huit le nombre de caméras en service : site A (caméras 2 et 3), site B (caméras 4 et 5), site C (caméra 6), site D (caméras 7 et 8) et site E (caméra 1). Les caméras 2, 3, 4, 5 et 6 ne sont ni retirées, ni déplacées, ni modifiées, et l'enregistreur continuera d'être exploité en mode standard d'enregistrement, les autres options restant désactivées.

L'efficacité du dispositif actuel, installé depuis peu, s'est d'ores et déjà vérifiée : le visionnage des enregistrements a permis de résoudre une incivilité grave, et d'assister les forces de l'ordre dans le cadre d'une investigation menée à la suite d'un cambriolage commis dans un village voisin.

L'extension a pour but d'inclure dans le dispositif des zones du village qui ne sont pas couvertes par les caméras actuelles, à savoir [zones non couvertes à préciser, correspondant aux sites D et E]. Le déplacement de la caméra 1, dont le champ de vision est aujourd'hui redondant avec celui de la caméra 2, permet par ailleurs de couvrir l'une de ces zones sans achat d'équipement supplémentaire.

Nous recevons à cet effet des remontées d'habitants concernés, déçus de ne pas être inclus dans le dispositif, et qui demandent à bénéficier de la même protection que le reste du village.

Ces demandes s'inscrivent dans le contexte persistant qui avait justifié l'installation initiale : [faits survenus depuis la mise en service : présence de véhicules suspects, vols, dégradations, constats des habitants].

Les finalités du dispositif demeurent inchangées, telles que définies à l'article 1er de l'arrêté précité :

\begin{itemize}
\item Protection des bâtiments et installations publics et de leurs abords
\item Prévention des atteintes à la sécurité des personnes et des biens dans des lieux particulièrement exposés à des risques d'agression et de vol ou de trafic de stupéfiant
\item Prévention et constatation des infractions relatives à l'abandon d'ordures, de déchets, de matériaux ou d'autres objets
\end{itemize}

Vous trouverez ci-joint le dossier de demande de modification, comprenant notamment le rapport de présentation, les plans de masse et de détail, la description du dispositif, la liste des personnes habilitées, les modalités d'information du public, le questionnaire de conformité de l'installateur, l'engagement de conformité et l'analyse d'impact relative à la protection des données.

Je me tiens à votre disposition et à celle de vos services, notamment du référent sûreté de la gendarmerie, pour tout complément utile à l'instruction de cette demande.

Je vous prie d'agréer, monsieur le Préfet, l'expression de mes meilleurs sentiments.

\closing{}
\end{letter}

\end{document}
